% This file is intended to be included from the main manuscript with \input.
% It uses the notation, theorem, assumption, algorithm, and algorithmic
% environments defined in the manuscript preamble.  All labels have the
% prefix ``spgd'' so that this section can coexist with earlier drafts.

\section{Soft-prefix Generalized DAPRO}
\label{sec:spgd}

This section gives a complete specification of the soft-prefix Generalized
DAPRO allocator used in our experiments.  The same allocation architecture is
used for LPB construction and population-level metric estimation.  The two
tasks differ only in the target event whose acquisition variance is placed in
the objective.  In particular, the score, the two-bin policy class, the
optimization algorithm, the projection to deployment rows, and the optional
CRC budget controller are shared.

The central idea is simple.  Hard Target-$A$ DAPRO puts the realized binary
target $A_i$ at the endpoint of each fully observed policy-fitting trajectory.
Soft-prefix Generalized DAPRO instead distributes this target over the
trajectory using the model's conditional event hazard at every observed
prefix.  This is a model-integrated, or Rao--Blackwellized, estimate of the
same target-specific inverse-propensity objective.  It is typically more
stable when the policy-fitting fold is small, while remaining a causal dynamic
policy because the decision at time $t$ uses only the prefix available at
time $t$.

\subsection{Sequential acquisition notation}

For prompt $i$, let $T_i$ be its one-based event time and let $Q_i$ be the
largest interaction that the acquisition policy may request.  We use
\begin{equation}
    Q_i=
    \begin{cases}
        \qprior(X_i)=\fhat_{\tauprior}(X_i),&\text{for LPB construction},\\
        \maxl,&\text{for metric estimation}.
    \end{cases}
    \label{eq:spgd_acquisition_horizon}
\end{equation}
The fully observed active length is
\begin{equation}
    L_i:=\min\{T_i,Q_i\}.
\end{equation}
Immediately before acquiring interaction $t$, the evaluator has observed the
current history $H_i(t)$ and chooses a conditional continuation probability
$P_i(t)\in[0,1]$.  Its cumulative reach probability is
\begin{equation}
    \rho_i(t):=\prod_{s=1}^{t}P_i(s),
    \qquad \rho_i(0):=1.
    \label{eq:spgd_cumulative_reach}
\end{equation}
Conditional on the complete trajectory, the expected acquisition cost of a
frozen policy is
\begin{equation}
    c_i(P):=\sum_{t=1}^{L_i}\rho_i(t).
    \label{eq:spgd_row_cost}
\end{equation}
The allocation is predictable: $P_i(t)$ may be a function of $X_i$,
$H_i(t)$, the training data, and the independent policy-fitting fold, but it
may not depend on $H_i(s)$ for $s>t$ or on the unobserved value of $T_i$.

\subsection{The task-specific target $A$}
\label{sec:spgd_task_target}

\paragraph{Metric estimation.}
For unsafe-event-rate estimation through the fixed horizon $\maxl$, the target
is
\begin{equation}
    A_i^{\mathrm{metric}}
    :=\ind{T_i\leq\maxl}.
    \label{eq:spgd_metric_target}
\end{equation}
The non-strict inequality is intentional: an event at interaction $\maxl$ is
part of the metric.

\paragraph{LPB construction.}
Let $\{\tau_j\}_{j=1}^{J}$ be the ordered candidate grid, restricted to
$\tau_j\leq\tauprior$.  For the raw-$\alpha$ implementation, freeze
\begin{equation}
    \tau_\star
    :=\max\{\tau_j:\tau_j<\tautarget\},
    \qquad \tautarget=0.10,
    \label{eq:spgd_raw_alpha_anchor}
\end{equation}
where the maximum is taken in the initial candidate prefix.  In continuous
notation one may write $\tau_\star=\tautarget$; the strict grid rule above is
the exact implementation.  Set
\begin{equation}
    F_i^\star:=\fhat_{\tau_\star}(X_i),
    \qquad
    A_i^{\mathrm{LPB}}
    :=\ind{T_i<F_i^\star}.
    \label{eq:spgd_lpb_target}
\end{equation}
The strict inequality agrees with the LPB miscoverage convention.  The anchor
is selected and frozen using only the policy-fitting fold before any Phase-II
acquisition outcomes are observed.

For either task, write $A_i$ for the appropriate target in
\eqref{eq:spgd_metric_target} or \eqref{eq:spgd_lpb_target}.  If $R_i$ is the
indicator that the acquisition process reaches the target event and
$\pi_i$ is its recorded inclusion probability, the HT contribution is
$R_iA_i/\pi_i$.  Conditional on the complete trajectories and a frozen
predictable policy,
\begin{equation}
    \operatorname{Var}\!\left(
      \frac{1}{N}\sum_{i=1}^{N}\frac{R_iA_i}{\pi_i}
      \middle|\mathcal D
    \right)
    =\frac{1}{N^2}\sum_{i=1}^{N}
      A_i\left(\frac{1}{\pi_i}-1\right).
    \label{eq:spgd_exact_ht_variance}
\end{equation}
Thus the target-specific design objective is $A$-weighted inverse reach, not
the unweighted mean inverse propensity.  For metric estimation,
\eqref{eq:spgd_exact_ht_variance} is exactly the acquisition variance of the
reported unsafe-event-rate estimator.  For an LPB, it is exact for the fixed
anchor $F_i^\star$; it is a proxy for the complete nonlinear procedure that
estimates an entire calibration curve and then selects a bound.

\subsection{Causal prefix hazards and soft target masses}

Let the separately trained prediction model return a conditional event-time
distribution after every available prefix.  The production score and soft
coefficient use its diagonal entry
\begin{equation}
    h_i(t)
    :=\widehat{\mathbb P}
      \bigl(T_i=t\mid T_i\geq t,H_i(t),\Dtrain\bigr).
    \label{eq:spgd_prefix_hazard}
\end{equation}
In the saved conditional probability tensor, this is
\texttt{conditional\_grid[i,t-1,t-1]}.  It is available before the draw that
decides whether interaction $t$ is acquired.  Using a later row of the tensor
would leak a future history and is not allowed.

Define the task mask
\begin{equation}
    M_i(t):=
    \begin{cases}
      \ind{t\leq\maxl},&\text{metric estimation},\\
      \ind{t<F_i^\star},&\text{LPB construction},
    \end{cases}
    \label{eq:spgd_target_mask}
\end{equation}
and the active-prefix mask $D_i(t):=\ind{t\leq L_i}$.  With the default global
regularization $\delta_{\mathrm{glob}}=0.001$, the exact soft event and cost
masses supplied to the optimizer are
\begin{align}
    a_i(t)
    &:=h_i(t)D_i(t)
      \frac{M_i(t)+\delta_{\mathrm{glob}}}
           {1+\delta_{\mathrm{glob}}},
      \label{eq:spgd_soft_mass}\\
    d_i(t)&:=D_i(t).
      \label{eq:spgd_cost_mass}
\end{align}
Equivalently, the numerator in \eqref{eq:spgd_soft_mass} mixes the named
target mass $h_i(t)D_i(t)M_i(t)$ with all event mass
$h_i(t)D_i(t)$ inside the acquisition envelope.  This small broad component
prevents target-empty directions from receiving exactly zero importance.  In
the metric task $Q_i=\maxl$ and $M_i(t)=1$ at every active prefix, so this
regularization cancels algebraically and $a_i(t)=h_i(t)D_i(t)$.

If the conditional hazards are correct, summing the soft masses over possible
event times integrates the hard target $A_i$.  The replacement therefore
reduces the label noise of putting one binary mass at the realized endpoint.
Correct hazard calibration is an efficiency condition, not a validity
condition: the final HT estimator remains design-unbiased under a
misspecified score model, provided that the deployed acquisition rule is
causal, its true propensities are recorded, and positivity holds.

The hazard has two distinct roles.  First, it determines the model-integrated
objective mass in \eqref{eq:spgd_soft_mass}.  Second, its value is the current
prefix score used to select a continuation probability.  These roles should
not be conflated: changing the target mask changes the objective without
changing the policy representation, while changing the score changes which
histories share a continuation probability.

\subsection{The hazard-K2 dynamic policy class}
\label{sec:spgd_k2}

Let $\calIone^{\mathrm{fit}}$ be the fully observed policy-fitting fold.  At
each interaction $t$, consider only rows with $L_i\geq t$ and compute the
empirical median
\begin{equation}
    \kappa_t
    :=\operatorname{Quantile}_{0.5}
       \{h_i(t):i\in\calIone^{\mathrm{fit}},\ L_i\geq t\}.
    \label{eq:spgd_median_cutpoint}
\end{equation}
The exact implementation removes duplicate cutpoints and assigns bins by
\begin{equation}
    b_i(t):=
    \begin{cases}
       0,&h_i(t)<\kappa_t,\\
       1,&h_i(t)\geq\kappa_t.
    \end{cases}
    \label{eq:spgd_bin_assignment}
\end{equation}
Thus a score tied at the median is assigned to the high bin.  If there is no
active policy-fitting row at time $t$, the policy uses structural probability
one in that column.  If one of the two cells is empty because of score ties,
its value is copied from the nearest observed cell.

The raw deployable policy has one parameter per time and score bin,
\begin{equation}
    P_i^{\mathrm{raw}}(t)=p_{t,b_i(t)},
    \qquad
    0<p_{t,0}\leq p_{t,1}\leq1.
    \label{eq:spgd_k2_policy}
\end{equation}
The monotonicity restriction expresses that a larger current event hazard
cannot receive a smaller continuation probability.  K2 therefore has at most
two \emph{raw current-score decisions} at a fixed time.  It does not reduce
the final allocation to two trajectory propensities.  The cumulative reach
\begin{equation}
    \rho_i^{\mathrm{raw}}(t)
      =\prod_{s=1}^{t}p_{s,b_i(s)}
\end{equation}
depends on the entire preceding sequence of low/high-bin assignments.  Two
rows in the same current bin can consequently have different cumulative
reach, and the cumulative projection below can also give them different
conditional continuation probabilities.  The method is therefore dynamic
and history-adaptive despite the deliberately low-capacity K2 response to the
new score at each individual time.

K2 is a regularization choice, not a mathematical requirement.  More bins or
a continuous monotone score map enlarge the policy class and can reduce the
population optimum, but they also estimate many more parameters from only
$N_1$ fully observed rows.  The production default uses K2 because it was the
most stable choice across small policy-fitting folds.

\subsection{The fitted convex objective}

Write
\begin{equation}
    y_{t,b}:=\log p_{t,b}\leq0,
    \qquad
    z_i(t):=\sum_{s=1}^{t}y_{s,b_i(s)},
\end{equation}
so $\rho_i^{\mathrm{raw}}(t)=\exp\{z_i(t)\}$.  For a policy-shape budget
$\bar B_{\mathrm{shape}}$, soft-prefix Generalized DAPRO solves
\begin{equation}
\begin{aligned}
    \min_{\{y_{t,b}\}}quad
    &\frac{1}{n_f}\sum_{i\in\calIone^{\mathrm{fit}}}
      \sum_{t=1}^{L_i}a_i(t)
      \left[\exp\{-z_i(t)\}-1\right]\\
    \text{subject to}\quad
    &\frac{1}{n_f}\sum_{i\in\calIone^{\mathrm{fit}}}
      \sum_{t=1}^{L_i}d_i(t)\exp\{z_i(t)\}
      \leq\bar B_{\mathrm{shape}},\\
    &y_{t,0}\leq y_{t,1}\leq0,
      \qquad t=1,\ldots,\maxl,
\end{aligned}
\label{eq:spgd_optimization}
\end{equation}
where $n_f=|\calIone^{\mathrm{fit}}|$.  The $-1$ term is constant in the
policy and may be omitted numerically.  Every objective term is an exponential
of an affine function of $y$, the budget is a convex exponential sublevel
constraint, and the score-monotonicity constraints are linear.  Hence
\eqref{eq:spgd_optimization} is a convex program for fixed cutpoints and bin
assignments.

The objective is the soft counterpart of the exact $A$-weighted variance in
\eqref{eq:spgd_exact_ht_variance}.  In the hard endpoint version, all row mass
is placed at $t=L_i$, giving
$A_i\{1/\rho_i(L_i)-1\}$.  In the soft version, the model distributes target
mass across possible event times and the optimizer protects each mass by the
cumulative reach needed to reveal it.

\subsection{Numerical optimization by a dual search and generalized PAV}

The implementation solves \eqref{eq:spgd_optimization} without fitting a
second regression from oracle probabilities to scores.  Introduce a
Lagrange multiplier $\lambda\geq0$ for the expected-cost constraint.  With all
other time columns held fixed, the coordinate loss at time $t$ has the form
\begin{equation}
    \sum_{i:L_i\geq t}
       \left\{\alpha_i(t)e^{-y_{t,b_i(t)}}
             +\beta_i(t)e^{y_{t,b_i(t)}}\right\},
    \label{eq:spgd_coordinate_loss}
\end{equation}
where
\begin{align}
    \alpha_i(t)
    &:=\frac{e^{y_{t,b_i(t)}}}{n_f}
      \sum_{u=t}^{L_i}a_i(u)e^{-z_i(u)},\\
    \beta_i(t)
    &:=\frac{\lambda e^{-y_{t,b_i(t)}}}{n_f}
      \sum_{u=t}^{L_i}d_i(u)e^{z_i(u)}.
\end{align}
For any adjacent score block $\mathcal B$ tied to a common log probability,
its unconstrained update is
\begin{equation}
    y_{\mathcal B}
    =\frac12\log
      \frac{\sum_{i\in\mathcal B}\alpha_i(t)}
           {\sum_{i\in\mathcal B}\beta_i(t)}.
    \label{eq:spgd_block_update}
\end{equation}
Generalized PAV merges adjacent score blocks whenever these updates violate
the nondecreasing score order, and clips the final values to
$[-700,0]$.  Because the input score is the bin identifier, all active rows
in the same time/bin cell are tied exactly.

For a fixed $\lambda$, the solver cycles over time coordinates and updates
the suffix cumulative log reaches after every PAV update.  It uses at most
$30$ inner sweeps and stops when the largest active log-probability change is
smaller than $5\times10^{-4}$.  An outer geometric bisection searches
$\lambda\in[10^{-12},10^{14}]$ for a budget-feasible solution, using at most
$60$ iterations and relative budget tolerance $10^{-9}$.  Among numerically
feasible candidates it retains the one with the smallest fitted objective,
breaking objective ties by distance to the budget boundary.

The resulting row-level probabilities are converted to a table by averaging
within every observed time/bin cell.  Empty cells copy the nearest observed
bin, and a cumulative maximum over bins restores
$p_{t,0}\leq p_{t,1}$.  Deployment uses this table directly.  Accordingly,
``projection'' in this version does not mean a learned Platt or isotonic
regression: it means evaluating the frozen empirical-median bin map and its
optimized lookup table on new prefix scores.

\subsection{Cumulative projection, positivity, and the no-CRC controller}
\label{sec:spgd_projection}

Let $r_i(t)$ be the cumulative product obtained from the frozen bin table.
One common cumulative-logit shift $\gamma$ is fitted on the policy-fitting
rows:
\begin{equation}
    \widetilde\rho_i(t;\gamma)
    :=\eps_\pi+(1-\eps_\pi)
      \operatorname{cummin}_{u\leq t}
      \sigma\!\left(\operatorname{logit}\{r_i(u)\}+\gamma\right),
    \label{eq:spgd_cumulative_projection}
\end{equation}
where $\sigma$ is the logistic function and
$\eps_\pi=0.005$ by default.  The cumulative minimum enforces temporal
non-increase, and the mixture ensures
$\widetilde\rho_i(t)\geq\eps_\pi$ on every active prefix.  A scalar bisection
chooses the largest $\gamma$ for which
\begin{equation}
    \frac{1}{n_f}\sum_{i\in\calIone^{\mathrm{fit}}}
      \sum_{t=1}^{L_i}\widetilde\rho_i(t;\gamma)
    \leq\bar B_{\mathrm{shape}}.
    \label{eq:spgd_projection_budget}
\end{equation}
Finally, cumulative reaches are converted back to conditional probabilities,
\begin{equation}
    \widetilde P_i(1)=\widetilde\rho_i(1),
    \qquad
    \widetilde P_i(t)=
      \frac{\widetilde\rho_i(t)}{\widetilde\rho_i(t-1)}.
    \label{eq:spgd_cumulative_to_conditional}
\end{equation}

Without CRC, all $N_1$ Phase-I rows are policy-fitting rows and are observed
fully through $L_i$.  If the total budget is $B=N\bar B$, define the remaining
Phase-II budget per row
\begin{equation}
    \bar B_{\mathrm{II}}
    :=\frac{N\bar B-\sum_{i\in\calIone}L_i}{N-N_1}.
\end{equation}
The production no-CRC method reserves the fixed projection-transfer margin
$\eta=1$ and uses
\begin{equation}
    \bar B_{\mathrm{shape}}:=\bar B_{\mathrm{II}}-\eta.
    \label{eq:spgd_no_crc_shape_budget}
\end{equation}
Its expected-total-budget claim therefore requires the transfer condition
that the deployment expected cost is at most the corrected policy-fitting
cost plus $\eta$.  The scalar bisection in
\eqref{eq:spgd_projection_budget} controls the empirical policy-fitting cost;
by itself it does not prove that the cost transfers to new trajectories.

\subsection{The optional independent CRC controller}
\label{sec:spgd_crc}

The CRC variant changes only budget control.  It does not change the target,
soft masses, hazard score, K2 policy class, or convex objective.  Split the
$N_1$ fully observed rows into a policy-fitting fold of size
$n_f=N_1-n_{\mathrm{crc}}$ and an independent control fold of size
$n_{\mathrm{crc}}$; in our experiments $n_{\mathrm{crc}}=N_1/2$.  The target
anchor, hazards, median cutpoints, bin table, logit shift, and candidate family
are all frozen before control-fold event times are inspected.  For fitting
the base shape, use
\begin{equation}
    \bar B_{\mathrm{shape}}
    :=\frac{N\bar B-\sum_{i\in\calIone^{\mathrm{fit}}}L_i}
            {N-n_f},
    \label{eq:spgd_crc_shape_budget}
\end{equation}
with no projection margin.

To obtain a useful bounded candidate-cost envelope, set
\begin{equation}
    C_{\max}:=\min\{\maxl,2\bar B\}.
    \label{eq:spgd_row_cap}
\end{equation}
Pool the fit-fold target masses,
$\bar a(t)=n_f^{-1}\sum_{i\in\calIone^{\mathrm{fit}}}a_i(t)$, and solve
\begin{equation}
\begin{aligned}
    \min_e\quad&\sum_{t=1}^{\maxl}\frac{\bar a(t)}{e(t)}\\
    \text{subject to}\quad
      &1\geq e(1)\geq\cdots\geq e(\maxl)\geq\eps_\pi,\\
      &\sum_{t=1}^{\maxl}e(t)\leq C_{\max}.
\end{aligned}
\label{eq:spgd_shared_envelope}
\end{equation}
Antitonic PAV on $\sqrt{\bar a(t)}$, followed by one scalar dual search,
solves \eqref{eq:spgd_shared_envelope}.  Intersect every base cumulative path
with this frozen shared envelope,
\begin{equation}
    \rho_i^{\mathrm{cap}}(t)
      :=\min\{\widetilde\rho_i(t),e(t)\}.
    \label{eq:spgd_apply_envelope}
\end{equation}
This operation is causal because both quantities are available at time $t$.
It also guarantees the pathwise support bound
$\sum_{t=1}^{Q_i}\rho_i^{\mathrm{cap}}(t)\leq C_{\max}$.  A row-specific
contraction computed from the row's complete future score path must not be
used; it would make an early decision depend on an unobserved future prefix.

Construct $K_{\mathrm{crc}}=401$ nested candidates with
$1=\xi_1>\cdots>\xi_{K_{\mathrm{crc}}}=0$ equally spaced:
\begin{equation}
    \rho_i^{(k)}(t)
    :=\eps_\pi+\xi_k
      \{\rho_i^{\mathrm{cap}}(t)-\eps_\pi\}.
    \label{eq:spgd_crc_candidates}
\end{equation}
For a fully observed control row, compute its pilot cost and the exact
conditional expected cost of candidate $k$,
\begin{equation}
    b_i:=L_i,
    \qquad
    c_i(k):=\sum_{t=1}^{L_i}\rho_i^{(k)}(t).
\end{equation}
Let $n=n_{\mathrm{crc}}$, $m=N-N_1$, and
$B_{\mathrm{rem}}=N\bar B-
\sum_{i\in\calIone^{\mathrm{fit}}}L_i$.  Define
\begin{equation}
    \ell_i(k):=c_i(k)+\frac{n}{m}b_i,
    \qquad
    L_{\max}:=C_{\max}+\frac{n}{m}\maxl.
\end{equation}
The CRC selector chooses the first, hence most aggressive, candidate satisfying
\begin{equation}
    \frac{\sum_{i\in\calIone^{\mathrm{crc}}}\ell_i(k)+L_{\max}}
         {n+1}
    \leq\frac{B_{\mathrm{rem}}}{m}.
    \label{eq:spgd_crc_rule}
\end{equation}
Under exchangeability of control and deployment rows, nesting, independence of
the control outcomes from the frozen family, and the stated cost bounds, this
gives marginal expected-total-budget control.  It is not a per-split or
realized-budget guarantee.  The factor $2$ in $C_{\max}=2\bar B$ is a practical
choice, not part of HT validity: $C_{\max}=\maxl$ is valid but creates a much
larger finite-sample CRC correction, while $C_{\max}=\bar B$ can prevent the
targeted policy from spending more than average on high-value rows.

\subsection{Deployment and recorded propensities}

For every Phase-II row, the evaluator repeats the following operation until
an event occurs, the row is censored, or $Q_i$ is reached.  At active time
$t$, it forms $H_i(t)$, evaluates the current hazard in
\eqref{eq:spgd_prefix_hazard}, applies the frozen median cutpoint, obtains the
base table value, applies the frozen cumulative projection and, when enabled,
the shared envelope and selected CRC contraction, converts cumulative reach
to a conditional probability, and draws the corresponding Bernoulli
continuation variable.  The exact product of the conditional probabilities
along the observed path is stored.  Fully observed policy-fit and control rows
have propensity one.

For the primary metric, an observed event at $T_i\leq\maxl$ has inclusion
propensity $\pi_i=\rho_i(T_i)$ and contributes $1/\pi_i$ to the HT total.  For
the fixed LPB candidate, an event satisfying $T_i<F_i^\star$ is treated
analogously.  The score model may be wrong without biasing this correction;
using the wrong propensity, allowing a Phase-II outcome to alter the frozen
policy, or allowing a future prefix to affect an earlier decision would break
the design argument.

\subsection{Should the LPB target use $\fhat_{\tautarget}$ or
$\fhat_{\tauprior}$?}
\label{sec:spgd_alpha_vs_prior}

Replacing the raw-$\alpha$ target by the prior bound is not a harmless
robustification; it changes the variance functional being optimized.  For an
ordered nested family, define
\begin{equation}
    A_i(\tau):=\ind{T_i<\fhat_\tau(X_i)},
    \qquad
    V_\tau(\rho):=\frac{1}{N^2}\sum_i
       A_i(\tau)\left(\frac{1}{\pi_i}-1\right).
\end{equation}
When $\tau\leq\tauprior$ and the candidates are nested,
$A_i(\tau)\leq A_i(\tauprior)$, and hence
\begin{equation}
    V_\tau(\rho)\leq V_{\tauprior}(\rho)
    \quad\text{for every fixed policy }\rho.
    \label{eq:spgd_prior_upper_bound}
\end{equation}
Thus the prior target controls a common, but potentially very loose, upper
bound on all smaller-candidate acquisition variances.  It does not minimize
each of them.  In particular,
\begin{equation}
    V_{\tauprior}(\rho)-V_{\tau_\star}(\rho)
    =\frac{1}{N^2}\sum_i
      \{A_i(\tauprior)-A_i(\tau_\star)\}
      \left(\frac{1}{\pi_i}-1\right).
    \label{eq:spgd_prior_band_cost}
\end{equation}
Optimizing the prior target therefore spends budget on every event in the
wide band between the raw-$\alpha$ candidate and the prior candidate.  When
$\tauprior$ is much larger than $\tautarget$---for example, $0.56$ versus
$0.10$---this can dilute the allocation precisely where the calibrated LPB
crossing is expected to occur.  There is no universal dominance theorem:
the policy minimizing $V_{\tauprior}$ can have substantially larger
$V_{\tau_\star}$ than the policy trained directly for $\tau_\star$.

For that reason, our default recommendation is to retain the raw-$\alpha$
target and its small broad regularizer in \eqref{eq:spgd_soft_mass}, rather
than replacing it by $\tauprior$.  The regularizer already protects all event
times through $Q_i$, including the prior endpoint, but gives them only weight
$\delta_{\mathrm{glob}}=0.001$ relative to the named target.

If substantial anchor misspecification is anticipated, a principled robust
extension is a weighted multi-target objective.  Choose a fit-fold-frozen grid
$\mathcal T_{\mathrm{rob}}\subseteq[0,\tauprior]$ and nonnegative weights
$w_\tau$ summing to one, concentrated on plausible values of the eventual
calibrated $\widehat\tau$.  Replace the LPB mask in
\eqref{eq:spgd_soft_mass} by
\begin{equation}
    M_i^{\mathrm{rob}}(t)
    :=\sum_{\tau\in\mathcal T_{\mathrm{rob}}}
       w_\tau\ind{t<\fhat_\tau(X_i)}.
    \label{eq:spgd_robust_mask}
\end{equation}
For hard targets this minimizes the weighted sum
\begin{equation}
    \sum_{\tau\in\mathcal T_{\mathrm{rob}}}w_\tau V_\tau(\rho)
    =\frac{1}{N^2}\sum_i
       \left\{\sum_\tau w_\tau A_i(\tau)\right\}
       \left(\frac{1}{\pi_i}-1\right),
    \label{eq:spgd_weighted_variances}
\end{equation}
so it remains exactly within the Generalized-DAPRO framework and can use the
same hazard-K2 optimizer.  A particularly conservative two-component choice
is
\begin{equation}
    M_i^{\mathrm{mix}}(t)
      =(1-\lambda)\ind{t<F_i^\star}
       +\lambda\ind{t<\fhat_{\tauprior}(X_i)},
    \qquad 0<\lambda\ll1.
    \label{eq:spgd_alpha_prior_mixture}
\end{equation}
Unlike wholesale replacement by $\tauprior$, this retains most effort near
the intended LPB boundary while explicitly purchasing robustness to a larger
calibration shift.  The weights, grid, and $\lambda$ must be selected using
training or policy-fitting information only, never the weighted Phase-II
calibration outcomes.

If the scientific target is genuinely worst-case accuracy over every
$\tau\in[0,\tauprior]$, one may instead solve the convex epigraph problem
$\min_\rho\max_{\tau\in\mathcal T}V_\tau(\rho)$.  That is a different robust
objective and requires a multi-constraint optimizer; setting
$A=A(\tauprior)$ is only a convenient upper-bound surrogate for it.

Finally, none of these fixed or weighted targets is exactly the variance of
the hard, data-selected LPB.  A selection-aware influence objective can in
principle target the calibration-curve crossing itself, but it must estimate
candidate switching sensitivities and is much less stable on small Phase-I
folds.  The raw-$\alpha$ soft-prefix objective is therefore the recommended
default: it is task-aligned, simple, causal, and empirically more stable than
the currently tested selection-aware alternatives.  A narrow weighted band
around $\tautarget$, rather than the entire interval up to $\tauprior$, is the
most defensible robustness ablation.

\subsection{Implementation checklist}

Algorithm~\ref{alg:spgd_exact} collects the complete implementation in one
place.

\begin{algorithm}[htbp]
\caption{Soft-prefix Generalized DAPRO with hazard-K2}
\label{alg:spgd_exact}
\begin{algorithmic}[1]
\REQUIRE Conditional prefix PMF grid; quantile grid $\{\fhat_{\tau_j}\}$;
$\tauprior$; total budget $N\bar B$; $N_1$; task type; optional CRC;
$\delta_{\mathrm{glob}}=0.001$; $\eps_\pi=0.005$; K2.
\STATE Set $Q_i$ by \eqref{eq:spgd_acquisition_horizon}; sample $N_1$ Phase-I
rows and observe each to $L_i=\min(T_i,Q_i)$.
\IF{CRC is enabled}
  \STATE Split Phase I into $n_f=N_1-n_{\mathrm{crc}}$ policy-fit rows and
  $n_{\mathrm{crc}}=N_1/2$ independent control rows.
\ELSE
  \STATE Use all $N_1$ rows for policy fitting.
\ENDIF
\STATE Freeze the task target: $T\leq\maxl$ for the metric, or the raw-$\alpha$
anchor in \eqref{eq:spgd_raw_alpha_anchor} for the LPB.
\FOR{each active policy-fit prefix $(i,t)$}
  \STATE Read $h_i(t)$ from the current-prefix diagonal PMF entry.
  \STATE Form $a_i(t)$ and $d_i(t)$ by
  \eqref{eq:spgd_soft_mass}--\eqref{eq:spgd_cost_mass}.
\ENDFOR
\FOR{$t=1,\ldots,\maxl$}
  \STATE On active policy-fit rows, compute the median hazard and assign bins
  by \eqref{eq:spgd_bin_assignment}; ties enter the high bin.
\ENDFOR
\STATE Set $\bar B_{\mathrm{shape}}$ by
\eqref{eq:spgd_no_crc_shape_budget} or \eqref{eq:spgd_crc_shape_budget}.
\STATE Solve \eqref{eq:spgd_optimization} by outer dual bisection and inner
generalized-PAV coordinate updates; construct the direct K2 lookup table.
\STATE Apply the shared cumulative-logit correction and propensity floor in
\eqref{eq:spgd_cumulative_projection}.
\IF{CRC is enabled}
  \STATE Fit the shared causal PAV envelope in
  \eqref{eq:spgd_shared_envelope}, form the $401$ candidates in
  \eqref{eq:spgd_crc_candidates}, and select one using only the control fold
  and \eqref{eq:spgd_crc_rule}.
\ENDIF
\STATE Freeze the target, cutpoints, table, cumulative shift, envelope, and
selected contraction before Phase-II outcomes are acquired.
\STATE Deploy causally.  At every reached prefix, evaluate the current hazard,
look up its K2 continuation, and record the exact product propensity.
\RETURN Censoring times and logged propensities for the LPB calibrator or HT
metric estimator.
\end{algorithmic}
\end{algorithm}

An independent implementation should additionally verify: one-based event
and horizon conventions; strict LPB and non-strict metric masks; exclusion of
inactive prefixes from both objective and cost; K2 cutpoints learned only on
active policy-fit rows; monotonicity of the two table entries; non-increasing
cumulative reach; the terminal floor; exact conversion between cumulative and
conditional probabilities; absence of future-prefix dependence; independence
of CRC outcomes from candidate construction; nesting and boundedness of CRC
candidate costs; and equality between every logged propensity and the product
of probabilities actually used by the sequential sampler.
